\chapter{O(N), Gross--Neveu, And Sigma-Model Families}
\label{ch:integrable-on-gross-neveu-sigma-model-families}

\ConventionsMixedCFT

This chapter introduces asymptotically free two-dimensional model families
whose exact scattering data carry nonabelian internal symmetry.  The aim is to
make explicit which pieces of the exact description are consequences of
symmetry, analyticity, and factorization, and which pieces require a
nonperturbative construction of the underlying local QFT.

\section{Internal-Symmetry Scattering Data}

Let \(G\) be a compact Lie group and let \(V\) be a finite-dimensional
unitary representation.  A massive particle multiplet with rapidity
\(\theta\) has one-particle space
\[
  \Hilb_1=L^2(\R,\dd\theta)\otimes V .
\]
A factorized two-particle scattering operator is a meromorphic family
\[
  S_{VV}(\theta):V\otimes V\longrightarrow V\otimes V,
\]
where \(\theta=\theta_1-\theta_2\).  \(G\)-invariance means
\[
  S_{VV}(\theta)\,(\rho(g)\otimes\rho(g))
  =
  (\rho(g)\otimes\rho(g))\,S_{VV}(\theta)
\]
for every \(g\in G\).  If
\[
  V\otimes V\simeq \bigoplus_\lambda m_\lambda V_\lambda ,
\]
then Schur decomposition gives
\[
  S_{VV}(\theta)
  =
  \bigoplus_\lambda S_\lambda(\theta)\otimes I_{V_\lambda}
\]
whenever the multiplicities \(m_\lambda\) are one.  With multiplicities, the
coefficient \(S_\lambda(\theta)\) is an endomorphism of the multiplicity
space.  Thus symmetry reduces the scattering problem to scalar or finite
matrix meromorphic functions, but unitarity, crossing, pole residues, and the
Yang--Baxter equation still constrain them.

\section{O(N)-Invariant Two-Particle Form}

For the vector representation \(V=\R^N\) of \(O(N)\), an invariant
two-particle amplitude has the index form
\[
  S_{ij}^{kl}(\theta)
  =
  \sigma_1(\theta)\delta_{ij}\delta^{kl}
  +\sigma_2(\theta)\delta_i^{\,k}\delta_j^{\,l}
  +\sigma_3(\theta)\delta_i^{\,l}\delta_j^{\,k}.
\]
Equivalently, decompose
\[
  V\otimes V
  =
  \mathbf 1\oplus \operatorname{Sym}^2_0(V)\oplus \wedge^2V .
\]
Let \(P_0,P_S,P_A\) be the corresponding projectors.  Then
\[
  S(\theta)=S_0(\theta)P_0+S_S(\theta)P_S+S_A(\theta)P_A .
\]
The two descriptions are related by applying the projectors to the index
formula.  For example,
\[
  P_0{}_{ij}^{kl}=\frac1N\delta_{ij}\delta^{kl},\qquad
  P_A{}_{ij}^{kl}=\frac12(\delta_i^{\,k}\delta_j^{\,l}
  -\delta_i^{\,l}\delta_j^{\,k}).
\]
The amplitude is unitary on the real rapidity axis when
\[
  S_R(\theta)S_R(-\theta)=1,\qquad
  R\in\{0,S,A\},
\]
after diagonalization in these channels.  Crossing identifies the particle
with its conjugate and relates \(S_{ij}^{kl}(\theta)\) to the amplitude at
\(\ii\pi-\theta\) with one external index crossed.

\section{Sigma-Model Ultraviolet Datum}

The Euclidean \(O(N)\) nonlinear sigma model is defined with field
\[
  n:\Sigma\longrightarrow S^{N-1}\subset\R^N,\qquad n\cdot n=1,
\]
and action
\[
  S_{\sigma}[n]
  =
  \frac1{2g^2}\int_\Sigma \dd^2x\,\partial_\mu n\cdot\partial_\mu n .
\]
The parameter \(g\) is the coefficient of the kinetic term in the displayed
regularized path integral.  Perturbation theory around a local coordinate
chart on \(S^{N-1}\) gives the one-loop running
\[
  \mu\frac{\dd g}{\dd\mu}
  =
  -\frac{N-2}{2\pi}g^3+O(g^5).
\]
The sign follows from the Ricci tensor of the round target:
\[
  R_{ab}=(N-2)h_{ab}.
\]
The corresponding sigma-model beta tensor has leading term
\[
  \mu\frac{\dd h_{ab}}{\dd\mu}
  =
  \frac1{2\pi}R_{ab}+O(R^2),
\]
and \(h_{ab}=g^{-2}h^{\mathrm{round}}_{ab}\) yields the displayed scalar
beta function.  A nonperturbative massive QFT construction must then provide
the Hilbert space, local fields or local algebras, and scattering states whose
two-particle data match the factorized \(O(N)\)-invariant amplitude.

\section{Gross--Neveu Ultraviolet Datum}

Let \(\psi^i\), \(i=1,\ldots,N\), be Majorana fermions in two Lorentzian
dimensions.  The \(O(N)\) Gross--Neveu action is
\[
  S_{\mathrm{GN}}
  =
  \int\dd^2x\,
  \left[
    \frac{\ii}{2}\psi^i\gamma^\mu\partial_\mu\psi^i
    +\frac{g^2}{8}(\bar\psi^i\psi^i)^2
  \right].
\]
Introducing an auxiliary scalar \(\sigma\) gives the equivalent finite-cutoff
functional
\[
  S_{\mathrm{GN}}[\psi,\sigma]
  =
  \int\dd^2x\,
  \left[
    \frac{\ii}{2}\psi^i\gamma^\mu\partial_\mu\psi^i
    -\frac12\sigma\,\bar\psi^i\psi^i
    -\frac1{2g^2}\sigma^2
  \right],
\]
because the Gaussian equation of motion is
\[
  \sigma=-\frac{g^2}{2}\bar\psi^i\psi^i .
\]
Substitution returns the quartic interaction.  The auxiliary-field form makes
the large-\(N\) mass gap equation explicit: for a constant saddle
\(\sigma=m\),
\[
  \frac{m}{g^2}
  =
  m\int^\Lambda\frac{\dd^2p_E}{(2\pi)^2}\frac1{p_E^2+m^2}.
\]
For \(m\ne0\),
\[
  \frac1{g^2}
  =
  \frac1{4\pi}\log\!\left(\frac{\Lambda^2+m^2}{m^2}\right).
\]
Thus a fixed physical mass \(m\) requires logarithmic dependence of \(g\) on
the regulator scale.  This derivation is a saddle calculation in a specified
regularized functional integral; the exact finite-\(N\) factorized QFT is a
separate construction.

\section{Bootstrap Matching Conditions}

For an asymptotically free model family, exact scattering data must match at
least four layers:
\[
  (G,V)\quad
  \longleftrightarrow\quad
  \text{particle multiplets},
\]
\[
  \beta(g)<0\quad
  \longleftrightarrow\quad
  \text{dynamically generated mass scale},
\]
\[
  \text{local currents and stress tensor}\quad
  \longleftrightarrow\quad
  \text{form-factor residues},
\]
\[
  \text{finite-volume spectrum}\quad
  \longleftrightarrow\quad
  \text{TBA kernels from }-\ii\partial_\theta\log S_R(\theta).
\]
The last line is checked by inserting the channel eigenvalues into the TBA
kernel
\[
  K_R(\theta)=\frac1{2\pi\ii}\partial_\theta\log S_R(\theta).
\]
The ultraviolet central charge extracted from the small-volume TBA must agree
with the short-distance fields that define the regulated model.  This is an
equality of two constructed objects only after the regulator-to-continuum map
and the scattering construction have been specified.

\begin{openproblem}[Constructive factorizing examples]
\label{op:constructive-factorizing-on-gn}
Construct the \(O(N)\) sigma model and Gross--Neveu factorizing QFTs as local
QFTs with Hilbert space, local algebras, conserved currents, mass gap,
Haag--Ruelle scattering states, and exact S-matrix, in one framework where the
ultraviolet regulator data and the bootstrap data are proved to match.
\end{openproblem}
